% Part: first-order-logic
% Chapter: natural-deduction
% Section: quantifier-rules

\documentclass[../../../include/open-logic-section]{subfiles}

\begin{document}

\olfileid{fol}{ntd}{qrl}

\olsection{పరిమాణీకరణిక నియమాలు}

\subsection{$\lforall$కు నియమాలు}

\begin{defish}
\AxiomC{$!A(a)$}
\RightLabel{\Intro{\lforall}}
\UnaryInfC{$\lforall[x][\Atom{!A}{x}]$}
\DisplayProof
\hfill
\AxiomC{$\lforall[x][\Atom{!A}{x}]$}
\RightLabel{\Elim{\lforall}}
\UnaryInfC{$!A(t)$}
\DisplayProof
\end{defish}

$\lforall$ నియమాల్లో $t$ ఒక సంవృత పదం (ఏ చరరాశులనూ కలిగి ఉండని
పదం); $a$ అనేది నిష్కర్ష~$\lforall[x][!A(x)]$లో గానీ,
పూర్వాధారం~$!A(a)$తో ముగిసే !!{derivation}లో !!{undischarged}గా ఉన్న
ఏ పరికల్పనలో గానీ కనిపించని ఒక !!a{constant}. \Intro{\lforall}
నిగమనం యొక్క \emph{ఐగెన్ చరరాశి} అని $a$ను అంటాం.\footnote{పై
నియమంలో $a$ ఒక స్థిరాంకమైనప్పటికీ “ఐగెన్ చరరాశి” అనే పేరునే
ఉపయోగిస్తాం. దీనికి చారిత్రక కారణాలున్నాయి.}

\subsection{$\lexists$కు నియమాలు}

\begin{defish}
\AxiomC{$\Atom{!A}{t}$}
\RightLabel{\Intro{\lexists}}
\UnaryInfC{$\lexists[x][\Atom{!A}{x}]$}
\DisplayProof
\hfill
\AxiomC{$\lexists[x][\Atom{!A}{x}]$}
\AxiomC{[$\Atom{!A}{a}$]$^n$}
\DeduceC{$!C$}
\DischargeRule{\Elim{\lexists}}{n}
\BinaryInfC{$!C$}
\DisplayProof
\end{defish}

ఇక్కడ కూడా $t$ ఒక సంవృత పదం; $a$ అనేది పూర్వాధారం
$\lexists[x][!A(x)]$లో గానీ, నిష్కర్ష~$!C$లో గానీ, రెండు
పూర్వాధారాలతో ముగిసే !!{derivation}sలో !!{undischarged}గా ఉన్న ఏ
పరికల్పనలో గానీ ($!A(a)$ పరికల్పనలు తప్ప) కనిపించని ఒక !!a{constant}.
\Elim{\lexists} నిగమనం యొక్క \emph{ఐగెన్ చరరాశి} అని $a$ను అంటాం.

\Intro{\lforall} లేదా \Elim{\lexists} నిగమనానికి దారితీసే
!!{derivation}sలోని పూర్వాధారాల్లో గానీ, !!{undischarged}గా ఉన్న ఏ
పరికల్పనలో గానీ ఐగెన్ చరరాశి కనిపించకూడదనే షరతును
\emph{ఐగెన్ చరరాశి షరతు} అంటాం.

\begin{explain}
!!{variable}~$x$ స్వేచ్ఛగా ఉన్న ఒక !!a{formula}ను $!A$ సూచిస్తే దాన్ని $!A(x)$గా
రాస్తామనే ఆచారాన్ని గుర్తుచేసుకోండి. అదే సందర్భంలో $!A(t)$ అనేది
$\Subst{!A}{t}{x}$కు సంక్షిప్త రూపం. కాబట్టి $\Intro\lexists$
నియమాన్ని ఇలా కూడా రాయవచ్చు:
\begin{prooftree}
\AxiomC{$\Subst{!A}{t}{x}$}
\RightLabel{\Intro{\lexists}}
\UnaryInfC{$\lexists[x][!A]$}
\end{prooftree}
$t$ ఇప్పటికే $!A$లో కనిపించవచ్చని గమనించండి; ఉదాహరణకు $!A$ అనేది
$\Atom{\Obj P}{t,x}$ కావచ్చు. అందువల్ల $\Atom{\Obj P}{t,t}$ నుంచి
$\lexists[x][\Atom{\Obj P}{t,x}]$ను నిగమించడం $\Intro\lexists$ను
సరిగ్గా వర్తింపజేయడమే---పూర్వాధారంలో $t$ కనిపించే స్థానాల్లో ఒకటి
లేదా అంతకంటే ఎక్కువ స్థానాలను బద్ధ !!{variable}~$x$తో “మార్చవచ్చు”;
అన్నిటినీ తప్పనిసరిగా మార్చనవసరం లేదు. అయితే $\Intro\lforall$,
$\Elim\lexists$లలోని ఐగెన్ చరరాశి షరతులు !!{constant}~$a$ అనేది
$!A$లో కనిపించకూడదని కోరుతాయి. కాబట్టి $\Atom{\Obj P}{a,a}$ నుంచి
$\lforall[x][\Atom{\Obj P}{a,x}]$ను $\Intro\lforall$ ఉపయోగించి
సరిగ్గా నిగమించలేం.
\end{explain}

\begin{explain}
\Intro{\lexists}, \Elim{\lforall}లలో ఎలాంటి పరిమితులూ లేవు; పదం~$t$
ఏదైనా కావచ్చు, కనుక షరతుల గురించి ఆందోళన అవసరం లేదు. మరోవైపు,
\Elim{\lexists}, \Intro{\lforall} నియమాల్లో ఐగెన్ చరరాశి షరతు,
!!{constant}~$a$ నిష్కర్షలో గానీ ఏ !!{undischarged} పరికల్పనలో గానీ
ఎక్కడా కనిపించకూడదని కోరుతుంది. వ్యవస్థ నిర్దుష్టంగా ఉండేలా, అంటే
ఏ !!{undischarged} పరికల్పనల నుంచి !!{sentence}s అనుగమిస్తాయో వాటి
నుంచే వాటిని !!{derive}s చేసేలా, ఈ షరతు అవసరం. ఈ షరతు లేకపోతే కింది
నిగమనం అనుమతించబడుతుంది:
\begin{prooftree}
  \AxiomC{$\lexists[x][!A(x)]$}
  \AxiomC{$\Discharge{!A(a)}{1}$}
  \RightLabel{*\Intro{\lforall}}
  \UnaryInfC{$\lforall[x][!A(x)]$}
  \RightLabel{\Elim{\lexists}}
  \BinaryInfC{$\lforall[x][!A(x)]$}
\end{prooftree}
అయితే $\lexists[x][!A(x)] \Entails/ \lforall[x][!A(x)]$.

పరిమాణీకరణికాల తొలగింపు నియమాలు !!{variable}s స్థానంలో సంవృత
పదాలను మాత్రమే ప్రతిస్థాపించనిస్తాయి. అందువల్ల !!{sentence}s సమితి
నుంచి నిగమించగల ప్రతి !!{formula} కూడా ఒక !!a{sentence}.
\end{explain}


\end{document}
